\documentclass[12pt]{article}

\usepackage[margin=1in]{geometry} % 1-inch margins
\usepackage{times}                % Times (Times New Roman–like) font
\usepackage{fancyhdr}             % For custom headers/footers
\usepackage{longtable}            % For multi-page tables
\usepackage{setspace}             % For 1.5 line spacing

\setcounter{page}{7}

\pagestyle{fancy}
\fancyhf{} % Clear default header and footer

% Header (Right-aligned)
\fancyhead[R]{\small Saksham AI: Multilingual Virtual Interviewer}

% Footer (Left-aligned: College Name, Right-aligned: Page Number)
\fancyfoot[L]{\small Pillai HOC College of Engineering and Technology (Autonomous), Rasayani}
\fancyfoot[R]{\small \thepage}

% Horizontal lines in header & footer
\renewcommand{\headrulewidth}{0.4pt} 
\renewcommand{\footrulewidth}{0.4pt} 

% Page numbering starts from 5
\setcounter{page}{5}

% Paragraph Formatting
\setlength{\parindent}{1.5em} % Indent first line
\setlength{\parskip}{0em}     % No extra space between paragraphs

\onehalfspacing  % 1.5 line spacing

\begin{document}

\vspace*{\fill}
\begin{center}
{\Large \bf Chapter 2} \\[2em]
{\Large \bf Literature Survey}
\end{center}
\vspace*{\fill}
\newpage

\setcounter{section}{2}

% Section 2.1
\subsection{Related Work}
\vspace{10pt}

\noindent The field of AI-based interview preparation and virtual recruitment systems has seen significant developments in recent years. Several platforms and technologies similar to “Saksham AI – Multilingual AI Interviewer” have been proposed and implemented to enhance candidate training and automated interview evaluation.
One notable related work is the AI Virtual Interview System (AIVIS) developed by Johnson et al. AIVIS uses Natural Language Processing and speech analysis to simulate real interview environments and evaluate candidate responses. The system includes features such as automated question generation, voice interaction, and performance feedback, which are similar to the core functionalities of Saksham AI.
Another relevant platform is the HireVue AI Interview Platform, which provides AI-driven video interview analysis, facial expression recognition, and communication skill assessment. HireVue focuses on improving recruitment efficiency by enabling automated candidate evaluation and structured interview processes using advanced artificial intelligence techniques.
In addition, the Mock Interview Training System (MITS) is another example of an AI-supported learning platform designed to help candidates practice interview questions and receive feedback on their answers. The system uses machine learning algorithms to analyze response quality, confidence level, and communication effectiveness.
These related works demonstrate the growing role of artificial intelligence in transforming interview preparation and recruitment processes, aligning with the objectives and features of “Saksham AI – Multilingual AI Interviewer”.


\newpage

% Section 2.2
\subsection{Literature Review}

\begin{longtable}{|c|p{4cm}|c|p{2.5cm}|p{2.5cm}|p{2.5cm}|}
\hline
\textbf{Sr. No.} & \textbf{Title of Paper} & \textbf{Year} & \textbf{Author(s)} & \textbf{Methodology} & \textbf{Limitations} \\
\hline
1 & AI-Powered Virtual Interview Assistants for Skill Development & 2025 & Dr. Vijayant Verma et al. & Proposes an AI-based multilingual interview system & Language translation inconsistency \\
\hline
2 & Multimodal Sentiment and Behavior Analysis in Virtual Interviews & 2025 & Anumeha Agrawal et al. & Uses deep learning (CNN + LSTM) for emotion recognition & Limited dataset diversity\\
\hline
3 & Language-Agnostic Interview Evaluation using NLP  & 2024 & Jyoti Sekhar Banerjee et al. & Employs BERT and RAG-based NLP models for cross-lingual question answering  & Limited conversational depth \\
\hline
4 & Virtual Recruitment Platforms with AI-Driven Assessment  & 2024 & Debashis De et al. & Presents an AI recruitment model with automatic scoring & The system lacks spoken communication analysis \\
\hline
5 & Real-Time Communication Skill Analysis through AI  & 2024 & Siddhartha Bhattacharyya et al. & Utilizes speech-to-text and sentiment analysis for evaluating clarity & The project is limited to single-language input \\
\hline
\end{longtable}

\begin{center}
Figure 2.2.1 Literature Review
\end{center}

\newpage

% Section 2.3
\subsection{Existing Systems}
\vspace{10pt}

In the current digital learning landscape, several platforms offer interview preparation and communication skill training through mock tests, question banks, or video tutorials. Popular systems such as InterviewBit, Pramp, and VMock provide structured interview practice, while platforms like Coursera, Udemy, and edX focus on theoretical and video-based soft skill development. However, these systems often function independently and lack integration between learning, practice, evaluation, and certification within a single framework.Most existing solutions focus on either knowledge enhancement or mock interviews, but not both in a connected and adaptive manner. Users frequently switch between different tools for interview tips, practice sessions, and feedback, resulting in fragmented and inefficient learning. Moreover, many of these systems do not provide real-time AI evaluation, multilingual support, or secure proctoring mechanisms, which are crucial for fair and realistic interview experiences.
\\

% Section 2.4
\subsection{Problem Statement}
\vspace{10pt}

Despite advancements in AI and communication technology, many candidates still struggle to develop effective interview skills and communicate confidently across languages. Traditional interview preparation methods and existing online platforms often lack interactive, personalized, and feedback-driven experiences that mirror real-world interview scenarios. This gap between theoretical preparation and practical performance limits learners’ ability to adapt to dynamic interview environments.Furthermore, most current systems focus only on question banks or recorded tutorials without providing real-time AI evaluation, speech analysis, or multilingual interaction. The absence of continuous feedback and performance tracking reduces engagement and progress consistency. With the increasing demand for AI-enabled, globally accessible interview training, there is a strong need for a solution that offers interactive practice, instant evaluation, and multilingual communication support—helping candidates build confidence, fluency, and professional readiness.
\end{document}